\documentclass[12pt,twoside,a4paper]{article}
\usepackage{amsmath}
\usepackage{datetime}
\usepackage[most]{tcolorbox}
\newcommand{\bash}[1]{\begin{tcolorbox}\texttt{\$ #1}\end{tcolorbox}}
\newcommand{\fn}[2]{\begin{tcolorbox}[enhanced,colback=white,colframe=gray!50!white,attach boxed title to top center={yshift=-3mm,yshifttext=-1mm},title={Funció \texttt{#1}},boxed title style={size=small,colframe=black},colbacktitle=white,coltitle=black]\texttt{#2}\end{tcolorbox}}

\title{Memòria Pràctica 1 - Mètodes Numèrics}
\author{Arnau del Río}
\date{\today}

\begin{document}
\maketitle

\section{Problema}
L'objectiu d'aquesta pràctica és poder treballar amb polinomis de grau elevat de forma àgil.

Les utilitats que desenvoluparem en aquesta pràctica són:
\begin{itemize}
    \item \texttt{pcl}: Calcular combinacions lineals de polinomis
    \item \texttt{pmul}: Multiplicar polinomis
    \item \texttt{pder}: Derivar polinomis
    \item \texttt{paval}: Avaluar polinomis
\end{itemize}

\newpage

\section{Implementació}
Per poder implementar aquestes utilitats, necessitarem poder guardar polinomis en memòria. Ho farèm a través de guardar un vector de coeficients, des de l'índex 0 (el terme independent) fins al grau del polinomi.

\subsection{Combinació lineal de polinomis}
El càlcul de una combinació lineal de polinomis consistirà en multiplicar cada coeficient del polinomi per el valor corresponent i sumar els resultats.

\begin{tcolorbox}[
title=\textbf{Algoritme:} Per calcular una combinació lineal $\alpha p(x) + \beta q(x)$,colback=red!5!white,colframe=red!75!black]
\begin{enumerate}
    \item Determinar el grau màxim $n$ entre els dos polinomis
    \item Per cada coeficient $i$ de 0 fins a $n$:
    \begin{itemize}
        \item Primer calculem el valor $r_i = \alpha p_i$
        \item Aleshores afegim $r_i = r_i + \beta q_i$
    \end{itemize}
\end{enumerate}
\end{tcolorbox}

\subsection{Multiplicació de polinomis}
La multiplicació de polinomis consistirà en multiplicar cada coeficient del polinomi per el valor corresponent i sumar els resultats.

\begin{tcolorbox}[
title=\textbf{Algoritme:} Per multiplicar dos polinomis $p(x)$ i $q(x)$,colback=blue!5!white,colframe=blue!75!black]
Siguin $n$ i $m$ els graus dels polinomis $p(x)$ i $q(x)$, respectivament. Aleshores, calculem:

$$r = \sum_{i = 0}^m \sum_{j = 0}^n p_i q_j x^{i+j}$$
\end{tcolorbox}

\subsection{Derivació de polinomis}
\begin{tcolorbox}[
title=\textbf{Algoritme:} Per derivar un polinomi $p(x)$,colback=green!5!white,colframe=green!75!black]
Determinar el grau $n$ del polinomi.

Per cada coeficient $i$ de 0 fins a $n-1$:

$$r_i = i \cdot p_i$$

Obtenim un nou polinomi $r(x)$ de grau $n-1$.
\end{tcolorbox}

\newpage

\section{Manual de software}
\subsection{Entrada i sortida dels polinomis}
Si volem passar un polinomi $p(x)$ de grau $n$ per \texttt{stdin}, passarem $n+2$ valors reals, on el primer valor és el grau del polinomi $n$ i els següents $n+1$ valors són els coeficients de $p(x)$.

Per exemple, per passar el polinomi $p(x) = 2x^3 - 3x^2 + x - 5$ a \texttt{pder}, passarem els valors $3$, $-5$, $1$, $-3$, $2$ per \texttt{stdin}:

\bash{echo 3 -5 1 -3 2 | ./pder}

La sortida per \texttt{stdout} tindrà el mateix format que l'entrada, en l'exemple anterior obtindrem la derivada del polinomi, $p'(x) = 6x^2 - 6x + 1$:

\bash{echo 3 -5 1 -3 2 | ./pder\\1 -6 6}

Les funcions \texttt{plleg} i \texttt{pescr} s'encarreguen de la lectura i escritura de polinomis per \texttt{stdin} i \texttt{stdout}, respectivament, i els seus arguments són un fitxer d'entrada o sortida (que pot ser \texttt{stdin} o \texttt{stdout}), un enter $n$ i un polinomi de grau $n$.

\subsection{Utilitat \texttt{pcl}}
La utilitat \texttt{pcl} es crida passant per \texttt{stdin} els valors de \texttt{alf} i \texttt{bet}, que són els coeficients dels polinomis a operar, i a continuació els polinomis a operar.
\bash{./pcl alf bet}

Exemple d'ús: si volem calcular $0.1 \cdot p(x) + 0.2 \cdot q(x)$, amb $p(x) = x^2 + x$ i $q(x) = 2x^2 - 3x + 1$, podem fer:

\bash{echo 2 1 0 1 2 -2 -3 1 | ./pcl 0.1 0.2}

\subsection{Funció \texttt{pcl}}
La funció \texttt{pcl} té el prototipus següent:

\fn{pcl}{void pcl (int m, double *p, int n, double *q, double alf, double bet, double r[]);}

\subsection{Utilitat \texttt{pmult}}
La utilitat \texttt{pmult} simplement multiplica dos polinomis passats per \texttt{stdin}.
\bash{./pmult p q}

Per exemple, per multiplicar els polinomis $p(x) = -1 + x$ i $q(x) = -2 + x$ amb resultat $r(x) = 2 - 3x + x^2$, podem fer:
\bash{echo 1 -1 1 1 -2 1 | ./pmult\\2 2 -3 1}

\subsection{Funció \texttt{pmult}}

\subsection{Utilitat \texttt{pder}}
La utilitat \texttt{pder} només reb un polinomi per \texttt{stdin} i escriu la seva derivada per \texttt{stdout}.
\bash{./pder p}

Per exemple, per derivar el polinomi $p(x) = 1 - 5x + 3x^2$ amb resultat $p'(x) = -5 + 6x$, podem fer:
\bash{echo 1 -5 3 | ./pder\\-5 6}

\subsection{Funció \texttt{pder}}

\subsection{Utilitat \texttt{paval}}
La utilitat \texttt{paval} evalua un polinomi passat per \texttt{stdin} en un conjunt de valors donats per \texttt{stdin} (sense necessitat de explicitar el nombre de valors).
\bash{./paval p x1 x2 ... xm }

Per exemple, per avaluar el polinomi $p(x) = -1 + x + x^2$ en els punts $x = 0.1, 0.2, 0.3$, podem fer:

\bash{echo 2 -1 1 1 0.1 0.2 0.3 | ./paval\\0.1 -0.89\\0.2 -0.76\\0.3 -0.61}

\subsection{Funció \texttt{paval}}

\newpage

\section{Validació del software}
Podem comprovar que

\newpage

\section{Resultats}
\subsection{gnuplot}

\end{document}
